\documentclass[11pt]{article}

\usepackage[letterpaper,margin=1in]{geometry}
\usepackage{amsmath,amssymb}
\usepackage{booktabs}
\usepackage{longtable}
\usepackage{array}
\usepackage{xcolor}
\usepackage{listings}
\usepackage{url}
\usepackage{hyperref}
\usepackage{titlesec}
\usepackage{parskip}
\usepackage{fancyhdr}
\usepackage{enumitem}

\hypersetup{
  colorlinks=true,
  linkcolor=black,
  urlcolor=blue!50!black,
  citecolor=black,
  pdftitle={KV Cache Engine x Precision Controller --- Integration Audit},
  pdfauthor={LonghornSilicon},
  pdfsubject={Integration Audit, Version kvce-acu-audit-0.2}
}

\titleformat{\section}{\normalfont\Large\bfseries}{\thesection.}{0.6em}{}
\titleformat{\subsection}{\normalfont\large\bfseries}{\thesubsection}{0.6em}{}
\titleformat{\subsubsection}{\normalfont\normalsize\bfseries}{\thesubsubsection}{0.6em}{}

\pagestyle{fancy}
\fancyhf{}
\lhead{\small KVCE $\times$ ACU Integration Audit}
\rhead{\small \texttt{kvce-acu-audit-0.2}}
\cfoot{\thepage}
\renewcommand{\headrulewidth}{0.4pt}

\definecolor{tabsep}{rgb}{0.93,0.93,0.93}

\setlist[itemize]{leftmargin=*,topsep=2pt,itemsep=1pt}

\title{\textbf{LonghornSilicon ACU $\times$ KV Cache Engine \\
       Integration Audit}}
\author{LonghornSilicon}
\date{Version \texttt{kvce-acu-audit-0.2} \quad 2026-05-22}

\begin{document}
\maketitle
\thispagestyle{fancy}

\begin{abstract}
\noindent
We audit the integrated dataflow between the Attention Compute Unit
(ACU, block~1) and the KV Cache Engine (KVCE, block~2) of the
LonghornSilicon LLM inference accelerator. The audit produced a
reference-model bug-fix (\texttt{kv-cache-engine@9b1163a}, closes
issue~\#1) that lifted decompressed-K cosine similarity from
$0.475$ to $0.978$ and reduced integrated relative MSE from $17.0$
to $0.36$ against a dense FP16 baseline. The chip-level integration
is unblocked. However, three quality gaps remain in the post-fix
state and are tracked here for traceability: a fixed-point boundary
clipping at the KVCE input port, a calibration mismatch at network
layer~0, and the absence of an end-to-end accuracy measurement on
real model outputs. Each is filed as a separate issue against the
KVCE repository.
\end{abstract}

\section{Scope and methodology}
\label{sec:scope}

The audit compares five attention-output paths, computed per~$64\!\times\!64$
attention tile on real Qwen2-0.5B traces (seq\_len $=512$, layers
$\{0,4,8,12,16,20,23\}$, $2744$ off-diagonal tiles total):

\begin{center}
\renewcommand{\arraystretch}{1.15}
\begin{tabular}{@{}llll@{}}
\toprule
\textbf{Path} & \textbf{K source} & \textbf{V source} & \textbf{S$\cdot$V matmul precision} \\
\midrule
A --- \textit{REF}              & true (FP16) & true (FP16) & FP16 dense \\
B --- PC alone (no KVCE)        & true (FP16) & true (FP16) & PC-routed \\
C --- KVCE alone, FP16 SV       & $\hat K$ (lossy) & $\hat V$ (lossy) & FP16 \\
D --- KVCE alone, INT8 SV       & $\hat K$ (lossy) & $\hat V$ (lossy) & INT8 \\
E --- Integrated                & $\hat K$ (lossy) & $\hat V$ (lossy) & PC-routed on lossy~$S$ \\
\bottomrule
\end{tabular}
\end{center}

Relative MSE of paths B, C, D, E against the REF baseline isolates
the contribution of each block: B measures the ACU in isolation;
C and~D measure the KVCE in isolation under the two possible routing
choices; E measures the integrated chip. Per-tile records carry
the precision controller's INT8/FP16 decision on both clean and
lossy~$S$, so we can also measure \emph{decision stability}
under upstream noise.

The harness lives at
\texttt{attention-compute-unit/analysis/integration\_test\_kv\_pc.py}
and runs in under two minutes per layer on an RTX~A4000. No tile,
attention matrix, or per-head activation is persisted to disk; per-tile
scalars accumulate in RAM and dump as a single JSON summary.

\section{Initial finding (pre-fix)}
\label{sec:initial}

The initial smoke pass (layer~0, $392$ tiles) returned median
relative MSE for the four non-REF paths:

\begin{center}
\begin{tabular}{@{}lr@{}}
\toprule
Path & Median rMSE \\
\midrule
B --- PC alone & $0.0001$ \\
C --- KVCE + FP16 SV & $17.01$ \\
D --- KVCE + INT8 SV & $17.05$ \\
E --- Integrated & $17.05$ \\
\bottomrule
\end{tabular}
\end{center}

The ACU's output error in isolation matched the original precision
controller paper's $91$--$97\%$ INT8-safe claim to four decimal
places. The KVCE-in-the-loop output error was $170{,}000\times$
larger, swamping the FP16-vs-INT8 distinction the precision
controller exists to make. Two isolation probes confirmed the
problem was upstream of the ACU entirely:

\begin{itemize}
\item On the KVCE repository's own \texttt{\_random\_vector} test
      inputs (which the repo's $184$ passing tests use), the median
      relative MSE of compress--decompress round-trip was $62\times$
      for keys and $1.9\times$ for values.
\item On charitable $\mathcal{N}(0, 1/\sqrt{d})$ Gaussian inputs
      (exactly the distribution the Lloyd-Max centroids were
      designed for), key reconstruction cosine similarity was
      $0.475$ --- a decompressed key vector pointed in essentially
      a random direction relative to the original.
\end{itemize}

The repository's existing round-trip test used an absolute-integer-space
MSE threshold of $\max_{\mathrm{val}}^2 \approx 10^9$ as a pass gate.
This is a \emph{doesn't-crash} threshold: literally random output
would pass it. There was no test asserting that reconstruction
approximated the input.

\section{Root cause and remediation}
\label{sec:rootcause}

Two bugs in the KVCE reference model, both in the fixed-point
arithmetic of the compress/decompress pipeline:

\begin{enumerate}
\item \textbf{Normalize shift bug.} In \texttt{kv\_cache\_engine\_ref.py},
the \texttt{normalize()} function divided the coordinate by the stored
norm using a left-shift of \texttt{coord\_frac} ($=12$) bits before
the integer division. The correct shift is \texttt{norm\_frac} ($=8$)
since the norm is stored in Q8.8 format. The $4$-bit overshift caused
a $16\times$ overscale on all post-normalize coordinates, which then
passed through the Walsh--Hadamard rotation and Lloyd--Max quantizer
in a regime far outside the centroid design range. After fix, key
cosine similarity rose from $0.475$ to $0.978$.

\item \textbf{C++ floor-division parity.} The C++ port used the
language's default \texttt{/} operator on negative numerators, which
truncates toward zero. The Python reference uses \texttt{//}, which
floors toward negative infinity. The two differ by one on negative
inputs and the per-coordinate error accumulated across the $64$-element
vector. A \texttt{floor\_div()} helper was added to the C++ port to
match the Python semantics bit-exactly.
\end{enumerate}

The fix is a three-line change in the Python reference and an
eight-line addition in the C++ port. The fix commit is
\texttt{kv-cache-engine@9b1163a} (closes
issue~\#1). The repository now ships a quality gate at
\texttt{sw/reference\_model/test\_reconstruction\_quality.py}
that asserts cosine similarity~$> 0.9$ and relative MSE~$< 0.1$ on
both K and V paths over Gaussian and \texttt{\_random\_vector}
input distributions.

\section{Post-fix results}
\label{sec:postfix}

\subsection{Aggregate}

Re-running the five-path integration harness on $2744$ off-diagonal
tiles across $7$ layers ($\{0,4,8,12,16,20,23\}$) yields:

\begin{center}
\renewcommand{\arraystretch}{1.15}
\begin{tabular}{@{}lrrr@{}}
\toprule
Path & Median rMSE & p90 & p99 \\
\midrule
B --- PC alone (no KVCE)    & $0.0002$ & $0.0014$ & $0.0250$ \\
C --- KV alone, FP16 SV     & $0.3637$ & $3.5709$ & $13.538$ \\
D --- KV alone, INT8 SV     & $0.3631$ & $3.5662$ & $13.447$ \\
E --- Integrated (PC + KV)  & $0.3631$ & $3.5662$ & $13.447$ \\
\bottomrule
\end{tabular}
\end{center}

The integrated median rMSE dropped from $17.0$ (pre-fix) to $0.36$
(post-fix), a $47\times$ improvement. The KVCE noise floor is now
$\sim 1800\times$ the precision controller's INT8 routing error,
down from $\sim 170{,}000\times$ pre-fix.

\subsection{Decision stability}

The precision controller's INT8/FP16 decision is computed on
$S = Q\hat K^\top / \sqrt{d}$ in the integrated path. Compared
with the decision computed on clean $S = Q K^\top / \sqrt{d}$:

\begin{center}
\renewcommand{\arraystretch}{1.15}
\begin{tabular}{@{}lr@{}}
\toprule
Metric & Value \\
\midrule
FP16\% on clean $S$ & $0.00\%$ \\
FP16\% on lossy $S$ & $0.18\%$ \\
Decisions agree & $99.82\%$ \\
Flipped clean-FP16 $\rightarrow$ lossy-INT8 & $0$ \\
Flipped clean-INT8 $\rightarrow$ lossy-FP16 & $5$ \\
\bottomrule
\end{tabular}
\end{center}

The decision function is robust to KVCE-induced perturbation of
the score matrix: $99.82\%$ of routing decisions are preserved,
and the five flipped tiles flipped \emph{conservatively}
(toward FP16, the safe direction).

\subsection{Per-layer breakdown}

\begin{center}
\renewcommand{\arraystretch}{1.15}
\setlength{\tabcolsep}{4pt}
\begin{tabular}{@{}rrrrrrrr@{}}
\toprule
Layer & $n$ & Clean FP16\% & Lossy FP16\% & rMSE B & rMSE C & rMSE D & rMSE E \\
\midrule
$0$  & $392$ & $0\%$ & $0\%$    & $0.0001$ & $2.567$ & $2.576$ & $2.576$ \\
$4$  & $392$ & $0\%$ & $0\%$    & $0.0002$ & $0.534$ & $0.534$ & $0.534$ \\
$8$  & $392$ & $0\%$ & $0\%$    & $0.0002$ & $0.254$ & $0.254$ & $0.254$ \\
$12$ & $392$ & $0\%$ & $0\%$    & $0.0003$ & $0.159$ & $0.160$ & $0.160$ \\
$16$ & $392$ & $0\%$ & $1.28\%$ & $0.0003$ & $0.335$ & $0.339$ & $0.339$ \\
$20$ & $392$ & $0\%$ & $0\%$    & $0.0004$ & $0.297$ & $0.304$ & $0.304$ \\
$23$ & $392$ & $0\%$ & $0\%$    & $0.0004$ & $0.342$ & $0.344$ & $0.344$ \\
\bottomrule
\end{tabular}
\end{center}

The $16\times$ spread between layer~$0$ ($2.567$) and layer~$12$
($0.159$) is the largest single per-layer variation and is the
subject of section~\ref{sec:gaps} item~\textit{(b)}.

\section{Residual quality gaps}
\label{sec:gaps}

The post-fix integrated rMSE of $0.36$ is $47\times$ better than
pre-fix but is still $\sim 3{,}600\times$ worse than the ACU in
isolation. Three distinct mechanisms contribute, ranked by
estimated magnitude:

\subsection*{(a) Q4.12 input clipping at the KVCE boundary}

The KVCE input port is specified as int16 Q4.12 fixed-point with
representable range $\pm 8.0$ at resolution $2^{-12}$. Real
Qwen2-0.5B post-\texttt{k\_proj} key activations have absolute
maximum approximately $\pm 152$ at the layers we measured. The
integration harness converts FP16~$\rightarrow$~Q4.12 at the
boundary, which clips approximately $95\%$ of $K$'s magnitude
information before any compression is attempted. Per-layer rMSE
correlates with $K$ magnitude, consistent with this being the
largest single contributor to the residual error. \textbf{This is
not a KVCE-internal defect}; it is a cross-block specification
gap. No document or issue currently records this constraint.
Filed as KVCE issue~\#2; see section~\ref{sec:openitems}.

\subsection*{(b) Layer-0 centroid mismatch}

The Lloyd--Max centroid set was derived under the assumption that
the Walsh--Hadamard-rotated normalized vector has coordinates drawn
from $\mathcal{N}(0, 1/\sqrt{d})$. This assumption is well approximated
for mid-network activations (layers $\geq 4$ in our trace) but
poorly approximated for layer~$0$, whose K/V outputs have not yet
been smoothed by the residual stream's accumulating LayerNorms.
The result is a $16\times$ rMSE spread across layers, with
layer~$0$ alone contributing the bulk of the corpus median.
Candidate remediations include per-layer centroid tables and an
adaptive \texttt{turbo4-adaptive} mode in the KVCE compression
algorithm. Filed as KVCE issue~\#3.

\subsection*{(c) TurboQuant+ \texttt{turbo4} compression noise floor}

The current KVCE configuration uses $3$-bit PolarQuant plus $1$-bit
QJL on keys ($4.25$~bpv) and $3$-bit PolarQuant only on values
($3.0$~bpv). At this compression ratio, per-vector reconstruction
error is approximately $0.04$ rMSE per coordinate. This error
compounds across the $64$~V-vectors of the $S \cdot V$ matmul,
producing the algorithmic floor that limits how good post-fix
integrated rMSE can be. Higher-precision modes ($6$~bpv
\texttt{turbo8}, $12$~bpv \texttt{turbo16}) would lower this floor
at the cost of compression ratio. This is a \emph{design property},
not a bug.

\section{Known open items}
\label{sec:openitems}

The following items are not blockers for cross-block dataflow but
are tracked for traceability through tape-out:

\begin{center}
\renewcommand{\arraystretch}{1.2}
\setlength{\tabcolsep}{5pt}
\begin{tabular}{@{}p{0.50\textwidth}lll@{}}
\toprule
\textbf{Item} & \textbf{Severity} & \textbf{Issue} & \textbf{Doc} \\
\midrule
Q4.12 input clipping vs. raw activation magnitudes
  & high   & kv-cache-engine\#2 & \S\ref{sec:gaps}(a) \\
Layer-0 centroid fit (per-layer or adaptive centroids)
  & medium & kv-cache-engine\#3 & \S\ref{sec:gaps}(b) \\
End-to-end accuracy on real LLM not yet measured
  & medium & kv-cache-engine\#4 & \S\ref{sec:gaps}, \S\ref{sec:impl-chip} \\
GQA convention in KVCE ISA
  & medium & --- (architectural)  & companion conflicts doc \\
\texttt{SRAM\_DEPTH} default is below hot-tier need
  & low    & --- (configuration)  & companion conflicts doc \\
\texttt{turbo4} noise floor dominates PC FP16 distinction
  & informational & --- (design property) & \S\ref{sec:gaps}(c) \\
\bottomrule
\end{tabular}
\end{center}

The first three items are tracked as KVCE GitHub issues with
this document referenced as the source. Items four and five are
covered in the companion document
\textit{Cross-Block Architectural Conflict Register}
(\texttt{kvce\_acu\_architectural\_conflicts.tex}).

\section{Interpretation}
\label{sec:interpretation}

\subsection*{Improvement from fix}

\begin{center}
\renewcommand{\arraystretch}{1.15}
\begin{tabular}{@{}lrrr@{}}
\toprule
Metric & Pre-fix & Post-fix & Improvement \\
\midrule
K cosine similarity                 & $0.475$    & $0.978$ & random $\rightarrow$ accurate \\
Integrated median rMSE              & $17.0$     & $0.36$  & $47\times$ \\
KVCE rMSE $/$ PC rMSE               & $170{,}000\times$ & $1{,}800\times$ & $94\times$ \\
\bottomrule
\end{tabular}
\end{center}

\subsection*{Answers to the original questions}

\textbf{Q1.\ Decision stability under lossy KV.}
Excellent. $99.82\%$ of precision-controller decisions are
unchanged when the input switches from true $K$ to lossy $\hat K$.
The five tiles that flipped all moved in the conservative
direction (toward FP16). The PC decision function is empirically
robust to $3$-bit PolarQuant plus $1$-bit QJL noise.

\textbf{Q2.\ Does FP16 routing still help under lossy~V?}
Not materially on the current $\texttt{turbo4}$ configuration.
The KV quantization rMSE ($0.36$) is roughly $1{,}800\times$ the
precision controller's routing-induced rMSE ($0.0002$), so the
FP16 vs.\ INT8 distinction ($0.0006$ rMSE between paths C and D)
is below the KV noise floor. The FP16 path is not actively harmful
--- the gate fires correctly --- but it cannot recover accuracy
that the upstream compression has already discarded. The
ACU's $1.7\times$ bandwidth-and-compute claim is intact for
\emph{ACU-isolated} workloads (e.g.\ external K/V supplied at
full precision); on the integrated chip, the claim degrades to
``bandwidth-and-compute savings with measurable accuracy cost
relative to dense FP16 attention.''

\textbf{Q3.\ Integrated MSE vs.\ ACU-alone MSE.}
Path E (integrated) tracks path D (KV+INT8) almost exactly,
because the PC routes virtually all tiles INT8 under lossy
conditions. The PC contributes neither benefit nor harm in this
configuration; it would re-acquire value at higher KVCE precision
modes where the K/V noise floor drops below the FP16 vs.\ INT8
distinction.

\subsection*{Layer-level observations}

\begin{itemize}
\item Layer~$0$ is the worst layer (rMSE $2.57$): early-layer
activations have not been LayerNorm-smoothed and do not match
the Lloyd--Max centroid design point.
\item Layer~$12$ is the best (rMSE $0.16$): mid-network activations
closest to the centroid distribution.
\item The $16\times$ spread across layers suggests per-layer or
adaptive centroid tables could substantially improve quality at
the same compression ratio.
\end{itemize}

\section{Implications per repository}
\label{sec:implications}

\subsection{\texttt{attention-compute-unit} (this repository)}

\begin{itemize}
\item The precision controller block is validated: rMSE $0.0002$
in isolation; decision stability $99.82\%$ under lossy inputs.
No code changes are required.
\item The FP16 routing path is not decorative in general --- it
remains valuable when V is not lossy-quantized (e.g.\ if a future
KVCE variant compresses K only, or under higher-precision modes).
On the current \texttt{turbo4} configuration, the KV noise
dominates and the PC FP16 path provides $< 0.2\%$ relative rMSE
improvement over INT8.
\item Future work: re-run this audit with higher-precision KVCE
modes (\texttt{turbo8}, \texttt{turbo16}) once they are implemented,
to test whether the FP16 distinction re-emerges as the KV noise
floor drops.
\end{itemize}

\subsection{\texttt{LonghornSilicon/kv-cache-engine}}

\begin{itemize}
\item Fix shipped: \texttt{kv-cache-engine@9b1163a}
(closes~\#1).
\item Quality gate added:
\texttt{sw/reference\_model/test\_reconstruction\_quality.py}
with cosine~$> 0.9$ and rMSE~$< 0.1$ thresholds. The previous
$\max_{\mathrm{val}}^2$ threshold did not bound reconstruction
quality.
\item Three follow-up issues filed (Q4.12 input range, Layer-0
centroid fit, end-to-end accuracy on real LLM) for the residual
quality gaps; see \S\ref{sec:gaps} and \S\ref{sec:openitems}.
\end{itemize}

\subsection{LonghornSilicon chip-level}
\label{sec:impl-chip}

\begin{itemize}
\item ACU integration is unblocked at the dataflow level: the two
blocks compose end-to-end with meaningful (non-noise) signal.
\item Chip-level accuracy claims (perplexity, downstream task
metrics on a real model) have not yet been measured against the
integrated dataflow. The per-tile rMSE of $0.36$ is suggestive
but does not directly bound model-output quality; this measurement
is tracked as kv-cache-engine\#4.
\item The binding constraint for end-to-end attention accuracy
on the current \texttt{turbo4} configuration is the KV cache
compression quality, not the SV matmul precision routing.
\end{itemize}

\section{Reproducing}
\label{sec:reproduction}

Verify the reference-model fix in isolation:

\begin{lstlisting}[language=Python,basicstyle=\ttfamily\small]
cd /home/<user>/kv-cache-engine
python3 sw/reference_model/test_reconstruction_quality.py
# Expect: K cosine ~ 0.978, V cosine ~ 0.985 on Gaussian inputs.
\end{lstlisting}

Re-run the integrated five-path harness:

\begin{lstlisting}[language=bash,basicstyle=\ttfamily\small]
cd /home/<user>/attention-compute-unit
python3 analysis/integration_test_kv_pc.py \
        --seq-len 512 --layers 0,4,8,12,16,20,23 \
        --prompts 1 --delete-cache-after
# Expect: integrated median rMSE ~ 0.36, layer-0 outlier ~ 2.57.
# Wall clock: ~2 min per layer on RTX A4000.
\end{lstlisting}

\section{Artifacts on this branch}
\label{sec:artifacts}

\begin{itemize}
\item \texttt{analysis/integration\_test\_kv\_pc.py} --- five-path
streaming harness, zero-tile-persistence, GQA-aware.
\item \texttt{analysis/integration\_test\_kv\_pc\_stats.json} ---
smoke-pass results for layer~$0$.
\item \texttt{docs/findings/kvce-acu-integration-audit.md} ---
prior markdown version of this audit, superseded by this LaTeX
document.
\item \texttt{docs/findings/kvce\_acu\_integration\_audit.tex} ---
this document.
\item \texttt{docs/findings/kvce\_acu\_architectural\_conflicts.tex} ---
companion document, cross-block conflict register covering items
beyond the immediate integration audit (GQA, SRAM depth, bridge
specification, etc.).
\end{itemize}

\section*{Change log}

\begin{itemize}
\item \texttt{kvce-acu-audit-0.1} (2026-05-19): initial markdown
audit; KVCE bug surfaced.
\item \texttt{kvce-acu-audit-0.2} (2026-05-22): post-fix re-run
($2744$ tiles, $7$ layers); promoted to LaTeX; residual quality
gaps catalogued; open items filed as KVCE issues.
\end{itemize}

\end{document}
